\documentclass[10pt]{article}

\usepackage[utf8]{inputenc}

\usepackage[T1]{fontenc}

\usepackage{amsmath}

\usepackage{amsfonts}

\usepackage{amssymb}

\usepackage[version=4]{mhchem}

\usepackage{stmaryrd}

\usepackage{mathrsfs}

\usepackage{graphicx}

\usepackage[export]{adjustbox}

\graphicspath{ {./images/} }



\begin{document}

Пусть функция $g(x)$ ограничена: $|g(x)| \leqslant c<\infty$. Тогда цена $s(x)$ является наименьшей эксцессивной мажорантой функции $g(x)$, т. е наименьшей из функций $f(x)$, удовлетворяющих двум свойствам



\[

g(x) \leqslant f(x), T f(x) \leqslant f(x),

\]



где $T f(x)=\mathrm{E}_{x} g\left(x_{1}\right)$. При этом момент



\[

\tau_{\varepsilon}=\inf \left\{n \geqslant 0: s\left(x_{n}\right) \leqslant g\left(x_{n}\right)+\varepsilon\right\}

\]



является $\varepsilon$-оптимальным для всякого $\varepsilon>0$, цена $s(x)$ удовлетворяет уравнению Вальда - Беллмана



\[

s(x)=\max \{g(x), T s(x)\}

\]



и может быть найдена по формуле



\[

s(x)=\lim _{n \rightarrow \infty} Q^{n} g(x),

\]



гдө



\[

Q g(x)=\max \{g(x), T g(x)\} .

\]



В том случае, когда множество $E$ конечно, момент



\[

\tau_{0}=\inf \left\{n \geqslant 0: s\left(x_{n}\right)=g\left(x_{n}\right)\right\}

\]



будет оптимальным. В общем случае момент $\tau_{0}$ является оптимальным, если $\mathrm{P}_{x}\left\{\tau_{0}<\infty\right\}=1, x \in E$.



Пусть



\[

C=\{x: s(x)>g(x)\}, \Gamma=\{x: s(x)=g(x)\} .

\]



В соответствии с определением



\[

\tau_{0}=\inf \left\{n \geqslant 0: x_{n} \in \Gamma\right\} .

\]



Иначе говоря, прекращение наблюдений следует производить при первом попадании в множество Г. В связи с этим множество $C$ наз. множеством продолжения наблюдений, а $\Gamma$ - множеством прекращения наблюдений.



Иллюстрацией этих результатов может служить задача различения двух простых гипотез, на к-рой А. Вальд продемонстрировал преимущество последовательных методов по сравнению с классическими. Пусть параметр $\theta$ принимает два значения 1 и 0 с априорными вероятностями $\pi$ и $1-\pi$ соответственно и множество заключительных решений $D$ состоит также из двух точек: $d=1$ (принимается гипотеза $H_{1}: \theta=1$ ) и $d=0$ (принимается гипотеза $\left.H_{0}: \theta=0\right)$. Если функцию $W_{1}(\theta, d)$ выбрать в виде



\[

W_{1}(\theta, d)=\left\{\begin{array}{l}

a, \theta=1, d=0, \\

b, \theta=0, d=1, \\

0, \text { в остальных случаях }

\end{array}\right.

\]



и положить



\[

W(\tau, \theta, d)=c \tau+W_{\mathbf{1}}(\theta, d),

\]



то для $R^{\delta}(\pi)$ получают выражение



\[

R^{\delta}(\pi)=c \mathrm{E}_{\pi} \tau+a \alpha_{\pi}(\delta)+b \beta_{\pi}(\delta),

\]



где



\[

\alpha_{\pi}(\delta)=\mathrm{P}_{\pi}\{d=0 \mid \theta=1\}, \beta_{\pi}(\delta)=\mathrm{P}_{\pi}\{d=1 \mid \theta=0\}

\]



\begin{itemize}

  \item вероятности ошибок первого и второго рода, а $P_{\boldsymbol{\pi}}$ оз- $^{-}$ начает распределение вероятностей в пространстве наблюдений, отвечающее априорному распределению $\pi$. Если $\pi_{n}=\mathrm{P}\left\{\theta=1 \mid \mathscr{F}_{n}\right\}-$ апостериорная вероятность гипотезы $H_{1} ; \theta=1$ относительно $\sigma$-алгебры $\mathscr{F}_{n}=\sigma\left\{\omega: \xi_{1}, \ldots, \xi_{n}\right\}$, то

\end{itemize}



\[

R^{\delta}(\pi)=E_{\pi}\left[c \tau+g\left(\pi_{\tau}\right)\right]

\]



где



\[

g(\pi)=\min (a \pi, b(1-\pi)) .

\]



Из общей теории оптимальных правил остановки, примененної к $x_{n}=\left(n, \pi_{n}\right)$, следует, что функция $\rho(\pi)=$ $=\inf _{\tau} R^{\delta}(\pi)$ удовлетворяет уравнению



\[

\rho(\pi)=\min \{g(\pi), c+T \rho(\pi)\} .

\]



Отсюда, в силу выпуклости вверх функций $\rho(\pi), g(\pi)$, $T \rho(\pi)$, можно вывести, что найдутся два числа $0 \leqslant$ $\leqslant A<B \leqslant 1$ такие, что область продолжения $C=\{\pi: A<$ $<\pi<B\}$, а область прекращения наблюдений $\Gamma=$ $=[0,1] \backslash(A, B)$. При әтом момент остановки



\[

\tau_{0}=\inf \left\{n \geqslant 0: \pi_{n} \in \Gamma\right\}

\]



является оптимальным $\left(\pi_{0}=\pi\right)$.



Если $p_{0}(x)$ и $p_{1}(x)$ - плотности распределений $F_{0}(x)$ и $F_{1}(x)$ (по мере $d \mu=\frac{1}{2}\left(d F_{0}+d F_{1}\right)$ ), а



\[

\varphi=\frac{p_{1}\left(\xi_{1}\right) \ldots p_{1}\left(\xi_{n}\right)}{p_{0}\left(\xi_{1}\right) \ldots p_{0}\left(\xi_{n}\right)}

\]



\begin{itemize}

  \item отношение правдолодобия, то область продолжения наблюдений (см. рис. 1) может быть записана в виде

\end{itemize}



\[

C=\left\{\varphi: \frac{A}{1-A} \frac{1-\pi}{\pi}<\varphi<\frac{B}{1-B} \frac{1-\pi}{\pi}\right\}

\]



и $\tau_{0}=\inf \left\{n \geqslant 0: \varphi_{n} \notin C\right\}$.



При этом если $\varphi_{\tau_{0}} \geqslant \frac{B}{1-B} \frac{1-\pi}{\pi}$, то выносится решение $d=1$, т. е. принимается гипотеза $H_{1}: \theta=1$. Если же $\varphi_{\tau_{0}} \leqslant \frac{A}{1-A} \frac{1-\pi}{\pi}$, то- гипотеза $H_{0}: \theta=0$.



Структура этого оптимального решающего правила сохраняется и для задачи различения гипотез в условноэкстремаль н о й постановке, состоящей в следующем. Для каждого решающего правила $\delta=(\tau$, d) вводят вероятности ошибок $\alpha(\delta)=P_{1}(d=$ $=0), \quad \beta(\delta)=P_{0}(d=1)$ и задают два числа $\alpha>0$ и $\beta>0$; и пусть, далее, $\Delta(\alpha, \beta)-$ сово-



\begin{center}

\includegraphics[max width=\textwidth]{2023_07_08_d993f03ee630042ae8b0g-1}

\end{center}



Рис. 1. купность всех решающих правил с $\alpha(\delta) \leqslant \alpha, \beta(\delta) \leqslant \beta$ и $\mathrm{E}_{0} \tau<\infty, \mathrm{E}_{1} \tau<\infty$. Следующий фундаментальный результат был получен А. Вальдом. Если $\alpha+\beta<1$ и среди критериев $\delta=(\tau$, $d)$, основанных на отношении правдоподобия $\varphi_{n}$ и имеющих вид



\[

\begin{gathered}

\tau=\inf \left\{n \geqslant 0: \varphi_{n} \notin(a, b)\right\}, \\

d=\left\{\begin{array}{l}

1, \varphi_{\tau} \geqslant b, \\

0, \varphi_{\tau} \leqslant a,

\end{array}\right.

\end{gathered}

\]



найдутся такие $a=a^{*}$ и $b=b^{*}$, что вероятности опибок первого и второго рода в точности равны $\alpha$ и $\beta$, то решающее правило $\delta^{*}=\left(\tau^{*}, d^{*}\right)$ с $a=a^{*}$ и $b=b^{*}$ является в классе $\Delta(\alpha, \beta)$ оптимальным в том смысле, что



\[

E_{0} \tau^{*} \leqslant E_{0} \tau, E_{1} \tau^{*} \leqslant E_{1} \tau

\]



для любого $\delta \in \Delta(\alpha, \beta)$.



Преимущества последовательного решающего правила $\delta^{*}=\left(\tau^{*}, d^{*}\right)$ по сравнению с классическим проще проиллюстрировать на примере задачи различения двух гипотез $H_{0}: \theta=0$ и $H_{1}: 0=1$ относительно локального среднего значения $\theta$ винеровского процесса $\xi_{t}$ с единичной диффузией. Оптимальное последовательное решающее правило $\delta^{*}=\left(\tau^{*}, d^{*}\right)$, обеспеqивающее заданные вероятности ошибок $\alpha$ и $\beta$ первого и второго рода соответственно, описывается следующим образом:



\[

\begin{aligned}

& \tau^{*}=\inf \left\{t \geqslant 0: \lambda_{t} \notin\left(a^{*}, b^{*}\right)\right\}, \\

& d^{*}=\left\{\begin{array}{l}

1, \lambda_{\tau *} \geqslant b^{*}, \\

0, \lambda_{\tau *} \leqslant a^{*},

\end{array}\right.

\end{aligned}

\]





\end{document}